% Appendix C -- attributed restatement of the decomposition theorem of the prior preprint.
\section{The weak-coupling decomposition theorem (prior work)}
\label{app:theorem}

This appendix restates, for completeness, the pathway-decomposition theorem that the present
work positions itself against. The theorem and its proof are the contribution of the earlier
BBCN118 preprint~\cite{bhatti2026} (CC BY 4.0); the statement and proof sketch below are
reproduced and adapted from that work. They are included here so that the boundary identified
in Section~\ref{sec:32} can be read without consulting the source. The new content is the final
remark, which names the assumption the present network violates.

\paragraph{Setup.} Let a Boolean network with global update $x(t{+}1)=F(x(t))$ be partitioned into
$m$ pathway modules $P_1,\dots,P_m$, with module states $x_i$ and module targets $x_i^{*}$, so that
$x=(x_1,\dots,x_m)$ and $x^{*}=(x_1^{*},\dots,x_m^{*})$.

\paragraph{Theorem 1 (convergence under pathway decomposition).} Assume:
\begin{itemize}
\item[(A1)] \emph{Weak inter-pathway coupling:} $\displaystyle\max_{i\neq j}\left\|\partial F_i/\partial x_j\right\| < \varepsilon$, with $\varepsilon \ll 1$;
\item[(A2)] \emph{Local kernel effectiveness:} for each module $i$ there is a kernel $u_i$ that drives $x_i(t)\to x_i^{*}$ in finitely many steps; and
\item[(A3)] \emph{Modularity:} $F_i(x)\approx f_i\!\left(x_i, e_i(x)\right)$, where $e_i(x)$ collects the inputs $P_i$ receives from other modules.
\end{itemize}
Then the sequential application of the pathway kernels $u=\bigcup_{i=1}^{m} u_i$ drives the global
state into a neighbourhood of the target,
\[
x(t)\to \mathcal{N}_\delta(x^{*})=\{x : d_H(x,x^{*})\le \delta\},\qquad \delta = O(\varepsilon),
\]
where $d_H$ is the Hamming distance.

\paragraph{Proof sketch.} Take the weighted Hamming Lyapunov function
\[
V(x)=\sum_{i=1}^{m} w_i\, d_H\!\left(x_i,x_i^{*}\right),\qquad w_i>0 .
\]
$V$ is non-negative, integer-valued, bounded above by $\sum_i w_i|P_i|$, and zero exactly at $x^{*}$.
Intervening on module $k$ reduces its term by some $\alpha_k>0$ by (A2). By (A1) every other module
moves by at most $w_j\varepsilon|P_j|$, so
\[
\Delta V \le -\alpha_k + \varepsilon\sum_{j\neq k} w_j|P_j| ,
\]
which is strictly negative once $\varepsilon$ is small enough relative to $\alpha_k$. Because $V$ is
bounded below, strictly decreasing, and integer-valued, it reaches a minimum $V_{\min}$ in finitely
many steps, with $V_{\min}\le\delta=O(\varepsilon)$; at $\varepsilon=0$ the bound gives $V_{\min}=0$
and exact convergence. \hfill$\square$

\paragraph{Corollary 1 (convergence rate).} If each pathway kernel reduces the mismatch by at least
$\beta>0$ per intervention, the system reaches the neighbourhood in at most $T \le V(x(0))\,/\,(\beta/2)$
interventions.

\paragraph{Remark: why this paper is in the other regime.} The result holds only under assumption
(A1). The present BBCN network does not satisfy it. Its pathway-coupling graph carries a dominant
strongly-connected component of 16 of the 22 regulatory pathways (Section~\ref{sec:32}), so the
inter-module coupling is not small, the $O(\varepsilon)$ neighbourhood is not tight, and pathway-wise
stabilisation does not deliver joint reachability. That boundary is precisely what motivates modelling
the apoptosis core directly as the bistable switch of Setup~B. The full proof, including the explicit
$\varepsilon$ threshold, is given in the cited preprint and is not reproduced here.